% Introduction.
%
% Order follows D19: the findings first, the suite third. The previous version
% led with the suite and with a method, and promised four metrics where three
% exist. Numbers here are quoted from the tables in Section 5.

\section{Introduction}
\label{sec:intro}

A world model is an agent's learned model of how its environment evolves. In
the decomposition of \citet{ha2018worldmodels} it has an encoder that maps
observations to latent states, a transition component $M$ that carries those
states forward under actions, and a controller trained inside the resulting
imagination. An agent that meets a sequence of environments must keep updating
$M$, and each update risks overwriting what $M$ knew about the environments
before --- catastrophic forgetting, at the component that every downstream use
of the model depends on. A transition model that has forgotten an environment's
dynamics generates incoherent rollouts for it no matter how good the policy is.

\citet{ha2018worldmodels} name this as an open problem and leave it to future
work. Since then it has been studied at the level of the agent: Continual World
\citep{wolczyk2021continualworld} measures forgetting in manipulation policies,
and \citet{kessler2023effectiveness} evaluate DreamerV2 as an integrated system
and report policy-level return, forgetting and transfer. Both measure what the
agent does. Neither isolates $M$, and we are not aware of a benchmark that
does. The distinction is not bookkeeping: a system-level result cannot say
whether a mitigation preserved the dynamics or preserved a policy robust to
their degradation, and those two ask for different fixes.

So we built the missing instrument --- three environment families, three levels
of dynamic distance each, five continual learning methods, ten seeds in every
cell that discriminates between them, 375 runs plus 75 from-scratch reference
pairs --- and this paper is mostly about what it told us that we had not asked. Both findings are negative, and neither is about
which method wins.

\paragraph{The distance axis does not order forgetting.} The benchmark, like
most in this literature, is built around an ordinal notion of how far apart two
tasks are, with the expectation that forgetting grows along it. It does not.
Forgetting peaks at the \emph{medium} level in all three families, so the
labelled order is wrong in every one of them. Across the nine cells the label
carries no rank information at all ($\rho = 0.00$), while both quantities we
measure --- the distance between the environments' transition models and the
difficulty of the new task --- do better ($0.58$ and $0.43$). We read the result
as forgetting scaling with what the new task demands of the model rather than
with its distance from the old one, on evidence set out in
Section~\ref{sec:results:predictors} that is firmer than either correlation.

The cleanest evidence needs no comparison across families. In Gymnasium the
maximum level is the medium perturbation plus two further ones, applied to the
same task pair; a strict superset of a perturbation cannot produce less
forgetting if the axis orders anything monotone, and it does. We rechecked that
cell at twice the training budget, because the obvious objection is that the
model simply never fits the harder task. The ordering holds, and the objection
turns out to be backwards: the more heavily perturbed task is \emph{easier} to
fit, since a cheetah at triple mass and $40\%$ gravity barely moves and
generates correspondingly little to learn. The strongest perturbation in that
family produces the poorest task, and the least forgetting.

\paragraph{The forgetting is in the encoder, where the metrics do not look.}
Component-level metrics have to hold a representation fixed in order to
attribute a change to the component, and ours do: prediction fidelity and
rollout divergence score $M$ in a latent basis encoded once. That is the
isolation we wanted, and it has a cost we can now measure. Fine-tuning's
held-out reconstruction of the first task degrades by a factor of $811$ while
its prediction fidelity reads $-5.58$ --- in the frozen basis, the model
\emph{improved}. EWC makes the dissociation exact and predicts it from its own
definition: its Fisher information is taken over $\log P(z' \mid z, a)$, which
contains no encoder parameter, so it is identically zero there. It preserves
the transition component almost perfectly ($-0.03$) while its encoder degrades
by $800$. A benchmark reporting only PF and RD would certify that neither model
forgot anything.

\paragraph{Contributions.}
\begin{enumerate}
  \item \textbf{A negative result about how these experiments are designed.}
  Level labels do not index forgetting, in any of three families, under any
  aggregation we tried, at two training budgets and at two seed counts. We recommend reporting a
  measured distance per task pair and holding it to account --- a standard our
  own measured distance fails, which we report rather than omit.
  \item \textbf{A negative result about how they are read.} Forgetting in a
  world model concentrates in the component that isolated metrics are blind to,
  and the most widely used regularizer cannot reach it for a structural reason
  that is checkable in advance. Component-level benchmarks owe their readers
  both scales.
  \item \textbf{The benchmark that produced them.} A component-level suite for
  forgetting in $M$: three metrics (prediction fidelity, rollout divergence,
  forward transfer), a measured distance between environments, three families,
  and a protocol that is declared in configuration, recorded in every result
  file, and refused by the runner if two budgets would meet in one cell. Five
  methods are characterised, not ranked.
\end{enumerate}

\paragraph{What this paper does not claim.} It does not claim that our own
method wins; it is one of the five, and the benchmark's clearest finding about
it is a failure mode (Section~\ref{sec:results:methods}). It does not report a
single aggregate score, because the one we started from is a median $88.5\%$
rollout divergence by weight. It does not report a policy-impact metric, which an
earlier version of this work announced and never implemented. And it supersedes
that version entirely: none of its five findings survive re-measurement, for
five instrumentation defects we document in Section~\ref{sec:discussion:previous},
including one that reverses. That history is part of the argument rather than
an embarrassment attached to it --- a benchmark's claim on anyone's attention is
that its numbers can be checked, and these were.
